\documentclass[11pt,a4paper]{article}
\usepackage[utf8]{inputenc}
\usepackage[T1]{fontenc}
\usepackage{lmodern}
\usepackage{amsmath,amssymb}
\usepackage{booktabs}
\usepackage{geometry}
\usepackage[colorlinks=true,linkcolor=blue!50!black,urlcolor=blue!50!black]{hyperref}
\usepackage{xcolor}
\geometry{a4paper,left=2.2cm,right=2.2cm,top=2.2cm,bottom=2.4cm}
\setlength{\parindent}{0pt}
\setlength{\parskip}{0.6ex}
\numberwithin{equation}{section}

\newcommand{\dg}{^{\dagger}}
\newcommand{\Dpv}{D_{\mathrm{PV}}}
\newcommand{\Cpv}{C_{\mathrm{PV}}}
\newcommand{\Spv}{S_{\mathrm{PV}}}
\newcommand{\Pp}{P_{+}}
\newcommand{\Pm}{P_{-}}
\newcommand{\Gfive}{\Gamma_5}
\newcommand{\AR}{\tilde A_{R}}
\newcommand{\AL}{\tilde A_{L}}
\newcommand{\Pro}{\mathcal P}
\newcommand{\Dc}{\mathcal D_{\mathrm c}}
\newcommand{\Dcpv}{\mathcal D_{\mathrm c}^{\mathrm{PV}}}
\newcommand{\Ce}{C_{e}}
\newcommand{\Co}{C_{o}}
\newcommand{\Cpve}{C_{\mathrm{PV},e}}
\newcommand{\Cpvo}{C_{\mathrm{PV},o}}
\newcommand{\code}[1]{\texttt{\small #1}}

\title{The factorized Pauli--Villars multigrid of RutticeQCD,\\ from the solver to HMC:
a step-by-step walkthrough}
\author{Wei-Lun Chen}
\date{\today\quad(internal note)}

\begin{document}
\maketitle

This is an internal walkthrough of the multigrid method for the M\"obius domain-wall operator as
it is implemented in RutticeQCD, written so that every formula can be followed line by line. It
starts from the operator, goes through the solver of $D\psi=\eta$, the setup of the multigrid and
the solves with $D\dg$, and ends with the solves of the HMC force. Nothing here is new with respect
to the MG note (\code{docs/mg\_background/MDWF\_MULTIGRID\_NOTE.tex}); the note states the results,
this walkthrough derives them step by step. The last section maps each step to the code.

\tableofcontents

% ======================================================================
\section{The M\"obius domain-wall operator}

\subsection{Ingredients}
The field $\psi$ lives on the four-dimensional lattice and carries a fifth index
$s=0,\dots,L_s-1$; each $\psi_s$ is a Dirac field. The chiral projectors are
\begin{equation}
P_\pm=\frac{1\pm\gamma_5}{2},\qquad P_\pm^2=P_\pm,\qquad P_+P_-=0,\qquad P_++P_-=1 .
\end{equation}
The Wilson hopping term and the Wilson operator of mass $-M_5$ are
\begin{align}
(K\psi)(x)&=\sum_{\mu=1}^{4}\Big[(1-\gamma_\mu)U_\mu(x)\psi(x+\hat\mu)
+(1+\gamma_\mu)U_\mu\dg(x-\hat\mu)\psi(x-\hat\mu)\Big],\\
W&=(4-M_5)-\tfrac12K .
\end{align}
$W$ acts on every slice $s$ alike and does not mix slices.

\subsection{The operator in Brower--Neff--Orginos form}
With $b$, $c$ the M\"obius parameters, define
\begin{equation}
D_+=1+bW,\qquad D_-=cW-1 .
\end{equation}
The operator at quark mass $m$ acts slice by slice as
\begin{equation}
(D\psi)_s=D_+\psi_s+D_-\big(\Pm\psi_{s+1}+\Pp\psi_{s-1}\big),
\qquad \psi_{L_s}\equiv-m\psi_0,\qquad \psi_{-1}\equiv-m\psi_{L_s-1}.
\label{eq:Dslice}
\end{equation}
For $L_s=4$ the matrix in $s$ is
\begin{equation}
D=\begin{pmatrix}
D_+ & D_-\Pm & 0 & -mD_-\Pp\\
D_-\Pp & D_+ & D_-\Pm & 0\\
0 & D_-\Pp & D_+ & D_-\Pm\\
-mD_-\Pm & 0 & D_-\Pp & D_+
\end{pmatrix}.
\end{equation}
The physical quark field is $q=\Pm\psi_0+\Pp\psi_{L_s-1}$; the runs use $b=1.5$, $c=0.5$,
$L_s=8$.

\subsection{Splitting off the fifth-dimension hop}
Define the hop in the fifth dimension, with a minus sign,
\begin{equation}
(Y\psi)_s=-\Pm\psi_{s+1}-\Pp\psi_{s-1}\qquad\text{(same wrap-around)}.
\end{equation}
Comparing with Eq.~\eqref{eq:Dslice}, the off-diagonal part of $D$ is $-D_-Y$, with $D_-$ on the
\emph{left}:
\begin{equation}
D=D_+-D_-Y .
\end{equation}
Now insert $D_\pm$ and then $W$, one substitution per line:
\begin{align}
D&=(1+bW)-(cW-1)Y\\
&=1+bW-cWY+Y\\
&=1+Y+W\,(b-cY)\\
&=1+Y+\Big[(4-M_5)-\tfrac12K\Big](b-cY)\\
&=\underbrace{\big[1+b(4-M_5)\big]}_{d_+}+\underbrace{\big[1-c(4-M_5)\big]}_{d_-}Y
-\tfrac12K\,(b-cY)\\
&=C-\tfrac12K\,(b-cY),\qquad C\equiv d_++d_-Y .
\label{eq:Dsplit}
\end{align}
$C$ contains no four-dimensional hop: it acts within one site, as an $L_s\times L_s$ matrix in
$s$. All the gauge field sits in $-\frac12K(b-cY)$. The Pauli--Villars operator is the same
operator at unit mass,
\begin{equation}
\Dpv=D\big|_{m=1},\qquad \Cpv=C\big|_{m=1}.
\end{equation}

\subsection{The site-diagonal part $C$ and its inverse}
$Y$ maps $\Pp\psi$ to $\Pp\psi$ and $\Pm\psi$ to $\Pm\psi$, so $C$ splits into a block on
$\psi^+=\Pp\psi$ and one on $\psi^-=\Pm\psi$. On $\psi^+$, $Y$ brings $\psi^+_{s-1}$ to slice $s$
with the factor $-1$, and across the wall the factor is $-(-m)=+m$:
\begin{equation}
(C\psi)^+_s=\begin{cases}
d_+\psi^+_0+m\,d_-\psi^+_{L_s-1}, & s=0,\\[2pt]
d_+\psi^+_s-d_-\psi^+_{s-1}, & s\ge1 .
\end{cases}
\end{equation}
To solve $C^+x=r$, take the rows $s\ge1$ first:
\begin{equation}
d_+x_s-d_-x_{s-1}=r_s\quad\Longrightarrow\quad x_s=\frac{r_s+d_-x_{s-1}}{d_+}.
\end{equation}
Each $x_s$ is then affine in the unknown $x_0$. With $q=d_-/d_+$ we claim
\begin{equation}
x_s=p_s+q^sx_0,\qquad p_0=0,\qquad p_s=\frac{r_s+d_-p_{s-1}}{d_+}.
\end{equation}
Check by induction:
\begin{align}
x_s&=\frac{r_s+d_-x_{s-1}}{d_+}\\
&=\frac{r_s+d_-(p_{s-1}+q^{s-1}x_0)}{d_+}\\
&=\frac{r_s+d_-p_{s-1}}{d_+}+\frac{d_-}{d_+}q^{s-1}x_0\\
&=p_s+q^sx_0 .
\end{align}
The remaining row $s=0$ fixes $x_0$:
\begin{align}
d_+x_0+m\,d_-x_{L_s-1}&=r_0\\
d_+x_0+m\,d_-\big(p_{L_s-1}+q^{L_s-1}x_0\big)&=r_0\\
\big(d_++m\,d_-q^{L_s-1}\big)x_0&=r_0-m\,d_-p_{L_s-1}\\
x_0&=\frac{r_0-m\,d_-p_{L_s-1}}{d_++m\,d_-q^{L_s-1}} .
\end{align}
The block on $\psi^-$ is the same with $s\to L_s-1-s$. So $C^{-1}$ costs $O(L_s)$ per site, is
exact, and does not depend on the gauge field. $C\dg$ has the transposed blocks and is inverted
the same way from the other end.

\subsection{Neff's $\alpha$}
Neff's operator is $D^{(\alpha)}=DQ$ with the site-local
\begin{equation}
Q_0=\Pm+\alpha\Pp,\qquad Q_s=\alpha\ (0<s<L_s-1),\qquad Q_{L_s-1}=\Pp+\alpha\Pm .
\end{equation}
$Q$ is Hermitian and $Q^{-1}$ is the same with $\alpha\to1/\alpha$. If $D^{(\alpha)}\psi^{(\alpha)}=\eta$, then
\begin{equation}
\eta=D^{(\alpha)}\psi^{(\alpha)}=DQ\psi^{(\alpha)}\quad\Longrightarrow\quad\psi=Q\psi^{(\alpha)}
\ \text{solves}\ D\psi=\eta .
\end{equation}
$Q$ leaves $\Pm\psi_0$ and $\Pp\psi_{L_s-1}$ alone, so the quark propagator is unchanged.

% ======================================================================
\section{The even-odd system}

\subsection{Blocks}
Order the sites even first. $K$ connects nearest neighbours, which have opposite parity, so in
Eq.~\eqref{eq:Dsplit} the diagonal blocks come from $C$ alone:
\begin{equation}
D=\begin{pmatrix}\Ce & D_{eo}\\ D_{oe} & \Co\end{pmatrix},\qquad
D_{pq}=-\tfrac12K_{pq}(b-cY)\quad(p\neq q).
\end{equation}
$\Ce$ and $\Co$ are the same matrix $C$ on the even and the odd sites.

\subsection{The Schur complement}
Solve $D\psi=\eta$ block by block. The odd rows:
\begin{align}
D_{oe}\psi_e+\Co\psi_o&=\eta_o\\
\psi_o&=\Co^{-1}\big(\eta_o-D_{oe}\psi_e\big).
\label{eq:psio}
\end{align}
Insert into the even rows:
\begin{align}
\Ce\psi_e+D_{eo}\psi_o&=\eta_e\\
\Ce\psi_e+D_{eo}\Co^{-1}\big(\eta_o-D_{oe}\psi_e\big)&=\eta_e\\
\big(\Ce-D_{eo}\Co^{-1}D_{oe}\big)\psi_e&=\eta_e-D_{eo}\Co^{-1}\eta_o\\
S\psi_e&=\eta'_e,\qquad S\equiv\Ce-D_{eo}\Co^{-1}D_{oe},\quad \eta'_e\equiv\eta_e-D_{eo}\Co^{-1}\eta_o .
\label{eq:schur}
\end{align}

\subsection{$S^{-1}=\pi D^{-1}\iota$}
Let $\iota$ put an even-site field on the whole lattice (zero on the odd sites) and $\pi$ take
the even part:
\begin{equation}
\iota y=\begin{pmatrix}y\\0\end{pmatrix},\qquad
\pi\begin{pmatrix}\psi_e\\\psi_o\end{pmatrix}=\psi_e .
\end{equation}
For $\eta=\iota\eta_e$ we have $\eta_o=0$, so $\eta'_e=\eta_e$ and Eq.~\eqref{eq:schur} reads
$S\psi_e=\eta_e$. Hence
\begin{equation}
S^{-1}\eta_e=\psi_e=\pi D^{-1}\iota\eta_e\quad\Longrightarrow\quad S^{-1}=\pi D^{-1}\iota .
\label{eq:Sinv}
\end{equation}
The same for $D\dg$: its block form is $\begin{pmatrix}\Ce\dg & D_{oe}\dg\\ D_{eo}\dg & \Co\dg\end{pmatrix}$,
and elimination gives $\Ce\dg-D_{oe}\dg(\Co\dg)^{-1}D_{eo}\dg=S\dg$, so $S^{-\dagger}=\pi D^{-\dagger}\iota$.

\subsection{The Pauli--Villars preconditioned forms}
$S$ is not suitable for a Krylov method: its spectrum surrounds the origin. Following Kanamori,
Chen and Matsufuru we multiply from the left by the adjoint of the Pauli--Villars factor and
divide the site part out. Two placements of $\Ce$ give two forms:
\begin{align}
\AR&=(\Spv\Cpve^{-1})\dg\,(S\Ce^{-1}),\\
\AL&=(\Cpve^{-1}\Spv)\dg\,(\Ce^{-1}S).
\end{align}
For $m=1$ both become $X\dg X$, Hermitian positive; for $m<1$ they are perturbations of it, with
all eigenvalues in the right half-plane.

\paragraph{Right form, step by step.} Start from $S\psi_e=\eta'_e$ and substitute $\psi_e=\Ce^{-1}y$:
\begin{align}
S\Ce^{-1}y&=\eta'_e\\
(\Spv\Cpve^{-1})\dg S\Ce^{-1}y&=(\Spv\Cpve^{-1})\dg\eta'_e\\
\AR\,y&=\rho_R,\qquad \rho_R\equiv(\Spv\Cpve^{-1})\dg\eta'_e,\\
\psi_e&=\Ce^{-1}y,\qquad \psi_o=\Co^{-1}(\eta_o-D_{oe}\psi_e).
\end{align}

\paragraph{Left form with $\alpha$ in $D$ only.} For $DQ$ the blocks are $\Ce Q_e$, $D_{eo}Q_o$,
$D_{oe}Q_e$, $\Co Q_o$, and the Schur complement is
\begin{align}
S^{(\alpha)}&=\Ce Q_e-D_{eo}Q_o(\Co Q_o)^{-1}D_{oe}Q_e\\
&=\Ce Q_e-D_{eo}Q_oQ_o^{-1}\Co^{-1}D_{oe}Q_e\\
&=\big(\Ce-D_{eo}\Co^{-1}D_{oe}\big)Q_e\\
&=SQ_e .
\end{align}
So the left Schur factor of $DQ$ is
\begin{equation}
K^{(\alpha)}=(\Ce Q_e)^{-1}SQ_e=Q_e^{-1}\Ce^{-1}SQ_e=Q_e^{-1}KQ_e,\qquad K\equiv\Ce^{-1}S,
\label{eq:Ksim}
\end{equation}
and the left form becomes $\AL^{(\alpha)}=(\Cpve^{-1}\Spv)\dg Q_e^{-1}(\Ce^{-1}S)Q_e$. In the right
form the two $Q_e$ cancel: $SQ_e(\Ce Q_e)^{-1}=S\Ce^{-1}$, so $\AR$ does not see $\alpha$.

% ======================================================================
\section{The outer solver: flexible GCR}
The outer iteration solves $\tilde Ay=\rho$ ($\tilde A=\AR$ or $\AL$) with a right preconditioner
$\mathcal M$, the multigrid cycle. Because $\mathcal M$ contains Krylov solves stopped by a
residual test, it changes from one application to the next, so the method must be flexible: it
keeps the preconditioned directions $z_i$ themselves.

One step of GCR($k$), from the residual $r_i=\rho-\tilde Ay_i$:
\begin{align}
z_i&=\mathcal M r_i &&\text{(one cycle)}\\
q_i&=\tilde Az_i &&\text{(one application of }\tilde A)\\
q_i&\leftarrow q_i-\sum_{j}\frac{(q_j,q_i)}{(q_j,q_j)}q_j,\quad
z_i\leftarrow z_i-\sum_{j}\frac{(q_j,q_i)}{(q_j,q_j)}z_j &&\text{(the last $k$ kept pairs)}\\
\alpha_i&=\frac{(q_i,r_i)}{(q_i,q_i)}\\
y_{i+1}&=y_i+\alpha_iz_i,\qquad r_{i+1}=r_i-\alpha_iq_i .
\end{align}
The coefficient $\alpha_i$ minimizes $\|r_i-\alpha q_i\|$, and the orthogonalization keeps
$\tilde Az_i=q_i$ true after the update, so $r_{i+1}=\rho-\tilde Ay_{i+1}$ exactly. With $k=2$
(GCR(2)) the code keeps $k+1=3$ pairs $(z,q)$ and the residual, seven fields on the even sites, at
the price of no optimality over the whole Krylov space.

% ======================================================================
\section{The two-grid cycle}
One cycle $y=\mathcal Mb$ is
\begin{align}
y&\leftarrow p(\tilde A)\,b &&\text{(pre-smoothing)}\\
r&\leftarrow b-\tilde Ay\\
y&\leftarrow y+e(r) &&\text{(coarse correction)}\\
r&\leftarrow b-\tilde Ay\\
y&\leftarrow y+p(\tilde A)\,r &&\text{(post-smoothing)}.
\end{align}

\subsection{The Chebyshev smoother, derived}
The smoother is $\nu$ steps of Chebyshev iteration on $\tilde Ax=b$ from $x=0$. Choose the interval
$[\lambda_{\rm lo},\lambda_{\rm hi}]$ that it damps ($\lambda_{\rm hi}$ slightly above the spectral
radius from a power iteration, $\lambda_{\rm lo}=\lambda_{\rm hi}/\kappa_s$) and set
\begin{equation}
\theta=\tfrac12(\lambda_{\rm hi}+\lambda_{\rm lo}),\qquad
\delta=\tfrac12(\lambda_{\rm hi}-\lambda_{\rm lo}),\qquad \sigma=\theta/\delta .
\end{equation}
The residual after $k$ steps is $r_k=R_k(\tilde A)r_0$ with the scaled Chebyshev polynomial
\begin{equation}
R_k(\lambda)=\frac{T_k\big((\theta-\lambda)/\delta\big)}{T_k(\sigma)},\qquad R_k(0)=1,
\end{equation}
which is at most $1/T_k(\sigma)$ in modulus on the interval. Write $t=(\theta-\lambda)/\delta$ and
use $T_{k+1}(t)=2tT_k(t)-T_{k-1}(t)$:
\begin{align}
T_{k+1}(\sigma)R_{k+1}&=2t\,T_k(\sigma)R_k-T_{k-1}(\sigma)R_{k-1}\\
R_{k+1}&=\rho_k\big[2t\,R_k-\rho_{k-1}R_{k-1}\big],\qquad \rho_k\equiv\frac{T_k(\sigma)}{T_{k+1}(\sigma)} .
\end{align}
The same recurrence at $t=\sigma$ gives the $\rho$'s:
\begin{align}
T_{k+1}(\sigma)&=2\sigma T_k(\sigma)-T_{k-1}(\sigma)\\
\frac1{\rho_k}&=2\sigma-\rho_{k-1},\qquad \rho_0=\frac{T_0(\sigma)}{T_1(\sigma)}=\frac1\sigma .
\end{align}
Now split $2t=2\sigma-2\lambda/\delta$ and use $2\sigma=1/\rho_k+\rho_{k-1}$:
\begin{align}
R_{k+1}&=\rho_k\Big[\big(2\sigma-\tfrac{2\lambda}{\delta}\big)R_k-\rho_{k-1}R_{k-1}\Big]\\
&=\rho_k\Big[\big(\tfrac1{\rho_k}+\rho_{k-1}\big)R_k-\tfrac{2\lambda}{\delta}R_k-\rho_{k-1}R_{k-1}\Big]\\
&=R_k+\rho_k\rho_{k-1}\big(R_k-R_{k-1}\big)-\frac{2\rho_k}{\delta}\lambda R_k .
\end{align}
Apply to $r_0$, with $\lambda\to\tilde A$:
\begin{equation}
r_{k+1}-r_k=\rho_k\rho_{k-1}(r_k-r_{k-1})-\frac{2\rho_k}{\delta}\tilde Ar_k .
\end{equation}
Since $r=b-\tilde Ax$, a step $x_{k+1}=x_k+d_k$ changes the residual by $-\tilde Ad_k$. Matching
the two sides,
\begin{equation}
d_k=\rho_k\rho_{k-1}\,d_{k-1}+\frac{2\rho_k}{\delta}\,r_k .
\end{equation}
The first step: $R_1=T_1(t)/T_1(\sigma)=t/\sigma=1-\lambda/\theta$, so $x_1=b/\theta$ and
$d_0=r_0/\theta$. These are the updates of the code (\code{BlasChebyshevStep}): $\nu$ steps cost
$\nu-1$ applications of $\tilde A$, and the polynomial is fixed, so the smoother is a linear map.
The runs use $\nu=8$, $\kappa_s=10$ with stout links ($M_5=1$) and $\nu=12$, $\kappa_s=30$ without
($M_5=1.8$).

\subsection{The factorized coarse correction}
\paragraph{The inverse of the right form.} Invert the product one factor at a time:
\begin{align}
\AR^{-1}&=\big[(\Spv\Cpve^{-1})\dg(S\Ce^{-1})\big]^{-1}\\
&=(S\Ce^{-1})^{-1}\,\big[(\Spv\Cpve^{-1})\dg\big]^{-1}\\
&=\Ce S^{-1}\,\big[(\Cpve^{-1})\dg\Spv\dg\big]^{-1}\\
&=\Ce S^{-1}\,(\Spv\dg)^{-1}\Cpve\dg\\
&=\Ce\;\pi D^{-1}\iota\;\;\pi(\Dpv\dg)^{-1}\iota\;\;\Cpve\dg ,
\label{eq:ARinv}
\end{align}
where Eq.~\eqref{eq:Sinv} was used twice. The inverse is a product of the inverses of two
\emph{one-hop} operators on the whole lattice, $D$ and $\Dpv\dg$, and site-local factors.

\paragraph{Galerkin replacement.} With a prolongation $\Pro$ (coarse to whole lattice), replace
each inverse by its Galerkin approximation:
\begin{equation}
X^{-1}\approx\Pro\big(\Pro\dg X\Pro\big)^{-1}\Pro\dg .
\end{equation}
Both coarse operators $\Dc=\Pro\dg D\Pro$ and $\Pro\dg\Dpv\dg\Pro$ are nearest-neighbour on the
coarse lattice, because $D$ is one hop and every aggregate contains the whole fifth dimension.
The correction of a residual $r$ (Eq.~\eqref{eq:ARinv} read from the right) is
\begin{equation}
e=O_L\,\pi\Pro\,\Dc^{-1}\Pro\dg\iota\;O_M\;\pi\Pro\big(\Pro\dg\Dpv\dg\Pro\big)^{-1}\Pro\dg\iota\;O_R\,r
\end{equation}
with the site-local factors
\begin{center}
\begin{tabular}{@{}lccc@{}}
\toprule
form & $O_R$ & $O_M$ & $O_L$\\
\midrule
right, any $\alpha$ & $\Cpve\dg$ & $1$ & $\Ce$\\
left, $\alpha$ in $D$ only & $1$ & $\Ce Q_e\Cpve\dg$ & $Q_e^{-1}$\\
left, $\alpha$ in $D$ and $\Dpv$ & $Q_e^{-1}$ & $\Ce Q_e^2\Cpve\dg$ & $Q_e^{-1}$\\
\bottomrule
\end{tabular}
\end{center}
In words: (i) put $O_Rr$ on the whole lattice, restrict, solve with the adjoint Pauli--Villars
coarse operator, prolong; (ii) take the even part, multiply by $O_M$, put it back on the whole
lattice, restrict, solve with the quark coarse operator, prolong; (iii) take the even part and
multiply by $O_L$. The coarse operators are always those of the plain $D$ and $\Dpv$, so one setup
serves both forms and every $\alpha$.

\paragraph{The left form with $\alpha$, step by step.} For $\alpha$ in $D$ only,
\begin{align}
\AL^{(\alpha)\,-1}&=\big[(\Cpve^{-1}\Spv)\dg\,Q_e^{-1}\Ce^{-1}SQ_e\big]^{-1}\\
&=\big(Q_e^{-1}\Ce^{-1}SQ_e\big)^{-1}\,\big[(\Cpve^{-1}\Spv)\dg\big]^{-1}\\
&=Q_e^{-1}S^{-1}\Ce Q_e\,\big[\Spv\dg(\Cpve^{-1})\dg\big]^{-1}\\
&=Q_e^{-1}S^{-1}\Ce Q_e\,\Cpve\dg\,(\Spv\dg)^{-1}\\
&=\underbrace{Q_e^{-1}}_{O_L}\;\pi D^{-1}\iota\;\underbrace{\Ce Q_e\Cpve\dg}_{O_M}\;
\pi(\Dpv\dg)^{-1}\iota\;\underbrace{1}_{O_R} .
\end{align}
Unlike the right form, the Galerkin factor of $\Dpv\dg$ acts on the plain residual, and
$\Cpve\dg$ lands between the two coarse solves, in $O_M$.

\subsection{Aggregation and the chirality split}
The fine lattice is divided into $4^4$ blocks, each with the whole fifth dimension; each block is
one coarse site. Every test vector $v$ is split by $\Gfive=\gamma_5R_5$, $(R_5\psi)_s=\psi_{L_s-1-s}$:
\begin{equation}
v_\pm=\tfrac12(1\pm\Gfive)v,\qquad \Gfive v_\pm=\pm v_\pm .
\end{equation}
With $n_v$ test vectors a coarse site has $2n_v$ components, the $+$ parts first. Each column of
$\Pro$ is then an eigenvector of $\Gfive$, so with $\Sigma=\mathrm{diag}(\mathbf 1_{n_v},-\mathbf
1_{n_v})$ per coarse site
\begin{equation}
\Gfive\Pro=\Pro\Sigma .
\end{equation}
\paragraph{Half the storage.} Write the condition $\Gfive v_\pm=\pm v_\pm$ out slice by slice:
\begin{align}
(\Gfive v_\pm)_s&=\gamma_5\,(R_5v_\pm)_s=\gamma_5\,v_\pm(L_s-1-s)\\
\pm v_\pm(s)&=\gamma_5\,v_\pm(L_s-1-s)\\
v_\pm(L_s-1-s)&=\pm\gamma_5\,v_\pm(s) .
\end{align}
So each part is fixed by its slices $s<L_s/2$; $\Pro$ is stored only there (``folded''), in half
precision. Restriction and prolongation unfold on the fly: for the restriction of a fine field $f$,
\begin{equation}
\big(\Pro\dg f\big)_{X,k}=\sum_{x\in X}\sum_{s<L_s/2}\Pro_{x,s;k}^{\,*}\,
\big(f_s\pm\gamma_5f_{L_s-1-s}\big)
\end{equation}
(the sign that of the chirality of column $k$), which is what \code{fold\_restrict\_tiled\_kernel}
computes.

\subsection{Orthonormalization per block: CholeskyQR2}
In each block and chirality the columns $w_1,\dots,w_n$ ($n=2\times$ the number of filtered
vectors, see below) are made orthonormal. Collect them as the columns of $W$ and write the QR
decomposition with a positive diagonal of $R$:
\begin{align}
W&=QR\\
W\dg W&=R\dg Q\dg QR=R\dg R\\
G\equiv W\dg W&=R\dg R\quad\text{(Cholesky)}\\
Q&=WR^{-1} .
\end{align}
So one Gram matrix ($n\times n$), one Cholesky factor and one triangular solve give $Q$; the
factorization is unique, so this $Q$ equals what modified Gram--Schmidt produces. In floating
point one pass loses orthogonality like the square of the condition number of $W$, so the pass
is done twice (CholeskyQR2): the second $G$ is close to the identity and repairs the first. Each
pass reads $W$ twice, against one sweep per pair of columns for Gram--Schmidt (800 sweeps for
$n=40$). Where a pivot of the first pass falls below $10^{-12}$ of its diagonal (nearly dependent
columns), the code falls back to Gram--Schmidt for those blocks.

\subsection{The basis $\{s,\hat t\}$ and why the correction is exact}
The setup produces vectors $y_j$ on the even sites (next section). Each is lifted to the whole
lattice in two ways, one for each coarse solve:
\begin{align}
s&=\big(\Ce^{-1}y,\;-\Co^{-1}D_{oe}\Ce^{-1}y\big),\\
\hat t&=\big(a,\;-(\Cpvo\dg)^{-1}D_{\mathrm{PV},eo}\dg a\big),\qquad a=\pi Ds .
\end{align}
\paragraph{$Ds$ vanishes on the odd sites.}
\begin{align}
Ds&=\begin{pmatrix}\Ce&D_{eo}\\D_{oe}&\Co\end{pmatrix}
\begin{pmatrix}\Ce^{-1}y\\-\Co^{-1}D_{oe}\Ce^{-1}y\end{pmatrix}\\
&=\begin{pmatrix}\Ce\Ce^{-1}y-D_{eo}\Co^{-1}D_{oe}\Ce^{-1}y\\
D_{oe}\Ce^{-1}y-\Co\Co^{-1}D_{oe}\Ce^{-1}y\end{pmatrix}\\
&=\begin{pmatrix}\big(\Ce-D_{eo}\Co^{-1}D_{oe}\big)\Ce^{-1}y\\0\end{pmatrix}\\
&=\iota\,S\Ce^{-1}y .
\end{align}
In the same way $\Dpv\dg\hat t=\iota\Spv\dg a$.
\paragraph{Exactness of the right form.} Let $s,\hat t\in\operatorname{range}\Pro$ and solve both
coarse systems exactly. The Galerkin map $\Pro(\Pro\dg X\Pro)^{-1}\Pro\dg X$ is the identity on
$\operatorname{range}\Pro$. Follow $r=\AR y$ through the correction ($O_R=\Cpve\dg$, $O_M=1$,
$O_L=\Ce$):
\begin{align}
O_Rr&=\Cpve\dg(\Spv\Cpve^{-1})\dg S\Ce^{-1}y\\
&=\Cpve\dg(\Cpve^{-1})\dg\Spv\dg S\Ce^{-1}y\\
&=\Spv\dg\,a,\qquad (a=S\Ce^{-1}y=\pi Ds)\\
\iota O_Rr&=\iota\Spv\dg a=\Dpv\dg\hat t\\
\Pro\big(\Pro\dg\Dpv\dg\Pro\big)^{-1}\Pro\dg\,\Dpv\dg\hat t&=\hat t
\qquad(\hat t\in\operatorname{range}\Pro)\\
\pi\hat t&=a=\pi Ds\\
\iota O_M\pi\hat t&=\iota\pi Ds=Ds\qquad(Ds\ \text{is zero on the odd sites})\\
\Pro\Dc^{-1}\Pro\dg Ds&=s\qquad(s\in\operatorname{range}\Pro)\\
e=O_L\pi s&=\Ce\Ce^{-1}y=y .
\end{align}
So the correction inverts $\AR$ exactly on the span of the $y_j$. In practice $\Pro$ is in half
precision and the coarse solves are inexact (relative residual $10^{-1}$), and the left form is
not exact at all; the cycle converges nonetheless.

\subsection{The coarse operators by probing}
$\Dc=\Pro\dg D\Pro$ has on each coarse site $X$ a $2n_v\times2n_v$ diagonal block $A(X)$ and one link
block $B_\mu^\pm(X)$ per direction: nine blocks. They are found column by column. For component $k$,
put a unit vector $\epsilon_k$ on \emph{every} coarse site and prolong:
$\phi=\Pro\epsilon_k$. Split $D=C+H$ ($H$ the hops) and $H$ by leg (direction and sign), and split
each leg into the part that stays inside its block and the part that crosses to the neighbour
block:
\begin{align}
D\phi&=C\phi+\sum_{\text{legs }l}H_l\phi\\
&=\underbrace{C\phi+\sum_lH_l^{\rm in}\phi}_{\text{stays in }X}
+\sum_l\underbrace{H_l^{\rm cross}\phi}_{\text{comes from }X\pm\hat\mu} .
\end{align}
Restricting, $\big(\Pro\dg(C\phi+\sum_lH^{\rm in}_l\phi)\big)_{X}$ is column $k$ of $A(X)$, and
$\big(\Pro\dg H^{\rm cross}_l\phi\big)_X$ is column $k$ of the link block of leg $l$, since every
coarse site carries the same unit vector and a crossing hop brings exactly the neighbour's
contribution. Doing this for all $2n_v$ components fills all nine blocks (\code{coarse\_blocks},
\code{hop\_split}, \code{restrict\_leg}). Several components are probed at a time, as many as a
quarter of the free device memory allows (four on $24^4$), so that $\Pro$ is read once for them.

\subsection{The coarse solves}
\paragraph{Quark factor.} $\Dc c=g$ by BiCGStab on its even-odd Schur complement, preconditioned by
the inverse diagonal blocks. With $\Dc=\begin{pmatrix}A_{ee}&B_{eo}\\B_{oe}&A_{oo}\end{pmatrix}$ on
the coarse lattice, exactly as in Eq.~\eqref{eq:schur},
\begin{align}
\big(A_{ee}-B_{eo}A_{oo}^{-1}B_{oe}\big)c_e&=g_e-B_{eo}A_{oo}^{-1}g_o,\\
c_o&=A_{oo}^{-1}\big(g_o-B_{oe}c_e\big).
\end{align}
\paragraph{Pauli--Villars factor.} $(\Pro\dg\Dpv\dg\Pro)w=g$. With $\Gfive\Pro=\Pro\Sigma$,
\begin{align}
\Pro\dg(\Gfive\Dpv\Gfive)\Pro&=\Pro\dg\Gfive\,\Dpv\,\Gfive\Pro\\
&=(\Gfive\Pro)\dg\,\Dpv\,(\Gfive\Pro)\\
&=\Sigma\Pro\dg\Dpv\Pro\Sigma=\Sigma\Dcpv\Sigma .
\end{align}
If $\Gfive\Dpv\Gfive=\Dpv\dg$ (true for $c=0$, off by $O(c)$ otherwise, Sec.~\ref{sec:dagger}),
the adjoint coarse factor is $\Sigma\Dcpv\Sigma$, and
\begin{align}
\Sigma\Dcpv\Sigma\,w&=g\\
\Dcpv\Sigma\,w&=\Sigma g\qquad(\Sigma^2=1)\\
\mathcal H_c\,w&=\Sigma g,\qquad \mathcal H_c\equiv\Dcpv\Sigma ,
\end{align}
$\mathcal H_c$ Hermitian (indefinite) for $c=0$, solved by MINRES with the absolute value of the
diagonal blocks as a positive preconditioner. Both factors stop at $|r|^2/|b|^2<10^{-2}$.

\paragraph{Levels below.} At small masses the coarse solves themselves become expensive (on
$24^4$ at $m=0.001$ about 2000 applications of $\Dc$ per correction). The coarse sites are then
grouped into super-blocks (2, then 3), local vectors are formed from $\Pro\dg s_j$, $\Pro\dg\hat
t_j$, and a V-cycle preconditions the coarse Krylov solve: with $Q_\ell$ the prolongation into
level $\ell-1$ and $B_\ell$ its diagonal blocks,
\begin{equation}
z=Q_{\ell+1}M_{\ell+1}Q_{\ell+1}\dg r,\qquad M_\ell r=z+B_\ell^{-1}\big(r-\mathcal D_\ell z\big),
\end{equation}
the last level inverted densely. This brings the coarse cost down to about 50 applications per
correction.

% ======================================================================
\section{Setup}
The setup produces the vectors $y_j$, then $\{s_j,\hat t_j\}$, $\Pro$, the coarse operators and the
levels. It always uses the right form.

\subsection{Chebyshev-filtered subspace iteration}
The $y_j$ should approximate the low right singular vectors of $\AR$, i.e.\ the low eigenvectors of
$H_n=\AR\dg\AR$. With $\ell(x)=\big(2x-(\lambda_{\rm hi}+\lambda_{\rm lo})\big)/(\lambda_{\rm
hi}-\lambda_{\rm lo})$, the polynomial $T_d(\ell(x))$ is bounded by $1$ on
$[\lambda_{\rm lo},\lambda_{\rm hi}]$ and grows fast below $\lambda_{\rm lo}$. One round on a block
$X$ of $m$ vectors:
\begin{align}
X&\leftarrow T_d\big(\ell(H_n)\big)X &&\text{(three-term recurrence, $d$ applications of $H_n$)}\\
X&\leftarrow\text{orthonormalize }X\\
G&=X\dg H_nX=V\Theta V\dg,\quad X\leftarrow XV &&\text{(Rayleigh--Ritz)}\\
\lambda_{\rm lo}&\leftarrow\theta_{\max} .
\end{align}
The degree is $d=1.5\sqrt{\lambda_{\rm hi}/\lambda_{\rm lo}}$, between 40 and 1000; the first round
uses $\lambda_{\rm lo}=\lambda_{\rm hi}/1000$. \emph{Observation} ($24^4$, $m=0.001$, 20 vectors): after
the first round (degree 48) the largest Ritz value is still $0.34\,\lambda_{\rm hi}$, so the
second round has degree 40 and damps only the top two thirds of the spectrum. The filter is weak;
the bootstrap does most of the work.

\subsection{Bootstrap}
Each round: build $\Pro$ and the cycle $\mathcal M$ from the current $y$; replace every $y_j$ by four
steps of FGMRES on $\AR x=y_j$ from $x=0$ preconditioned by $\mathcal M$; orthonormalize $y$ through
its Gram matrix. Why it works: write $y_j$ in eigenvectors of the preconditioned operator. The
cycle already handles the components inside its coarse space, and FGMRES with it removes them
first; what remains, and is therefore amplified relative to the rest, is what the coarse space
misses. In HMC the levels below are used from the second round on (in the first, on the filtered
vectors, the coarse solves were as fast without them).

\subsection{Cost and memory}
On $24^4$ ($L_s=8$, 20 filtered vectors, $m=0.001$, A30): the filter about 35\,s and the bootstrap
about 77\,s, the setup in all about 113\,s. A bootstrap round makes its correction ($\Pro$ 1.8\,s,
the two coarse operators about 4.4\,s each, the inverse blocks 0.8\,s, the top of the smoother
1.3\,s; the levels 7\,s more in the second round) and then improves the 20 vectors (about 20\,s in
the first round, 26\,s in the second). $\Pro$ ($40$ columns, half precision, whole lattice) is the largest
object, about 4.7\,GB.

% ======================================================================
\section{Solving with $D\dg$}
\label{sec:dagger}
HMC needs $(D\dg D)^{-1}$, so it solves with $D\dg$ as often as with $D$. A second setup for $D\dg$
would double the cost. The tool that avoids it is an exact relation between $D$ and $D\dg$.

\subsection{$\Gfive D\Gfive\neq D\dg$}
Each building block transforms simply under $\Gfive=\gamma_5R_5$:
\begin{align}
\Gfive W\Gfive&=W\dg,\qquad \Gfive D_\pm\Gfive=D_\pm\dg,\qquad \Gfive Y\Gfive=Y\dg .
\end{align}
Then, from $D=D_+-D_-Y$,
\begin{align}
\Gfive D\Gfive&=\Gfive D_+\Gfive-\Gfive D_-\Gfive\,\Gfive Y\Gfive=D_+\dg-D_-\dg Y\dg,\\
D\dg&=D_+\dg-(D_-Y)\dg=D_+\dg-Y\dg D_-\dg .
\end{align}
The order of the last product differs. Subtracting,
\begin{align}
\Gfive D\Gfive-D\dg&=Y\dg D_-\dg-D_-\dg Y\dg\\
&=Y\dg(cW\dg-1)-(cW\dg-1)Y\dg\\
&=c\,\big(Y\dg W\dg-W\dg Y\dg\big)=c\,[Y\dg,W\dg],
\end{align}
which is of order $c$ (zero for Shamir, $0.15$ relative on a $12^4$ stout configuration at
$c=0.5$).

\subsection{The exact relation}
$D_+$ and $D_-$ are functions of $W$, so they commute. Factor $D_-$ out to the left (formally):
\begin{equation}
D=D_-K_D,\qquad K_D=D_-^{-1}D_+-Y .
\end{equation}
$K_D$ is $\Gfive$-Hermitian:
\begin{align}
\Gfive K_D\Gfive&=\Gfive D_-^{-1}\Gfive\,\Gfive D_+\Gfive-\Gfive Y\Gfive\\
&=(D_-\dg)^{-1}D_+\dg-Y\dg\\
&=D_+\dg(D_-\dg)^{-1}-Y\dg\qquad(D_\pm\dg\text{ commute})\\
&=(D_-^{-1}D_+)\dg-Y\dg=K_D\dg .
\end{align}
So
\begin{align}
D\dg&=K_D\dg D_-\dg\\
&=\Gfive K_D\Gfive\,D_-\dg\\
&=\Gfive D_-^{-1}D\,\Gfive\,D_-\dg ,
\end{align}
and multiplying by $D_-\Gfive$ from the left,
\begin{equation}
\boxed{\;D_-\,\Gfive\,D\dg=D\,\Gfive\,D_-\dg\;}
\end{equation}
with no inverse left. $D_-$ itself cannot be inverted in practice: $D_-=c(W-1/c)$ is $c$ times a
Wilson operator of mass $-M_5-1/c=-3$, whose spectrum surrounds the origin (GCR stalls at
$|r|^2/|b|^2\approx0.74$).

\subsection{$\Gfive D^{-1}\Gfive$ as a preconditioner of $D\dg$}
Multiply $D\dg$ from the right by $\Gfive D^{-1}\Gfive$ and use the relation above in the form
$D\dg=\Gfive D_-^{-1}D\,\Gfive D_-\dg$, together with $\Gfive D_-\dg\Gfive=D_-$:
\begin{align}
D\dg\,\Gfive D^{-1}\Gfive&=\Gfive D_-^{-1}D\,\Gfive D_-\dg\,\Gfive\,D^{-1}\Gfive\\
&=\Gfive D_-^{-1}D\,D_-\,D^{-1}\,\Gfive\\
&=\Gfive D_-^{-1}\big(D_-D+[D,D_-]\big)D^{-1}\Gfive\\
&=\Gfive\big(1+D_-^{-1}[D,D_-]D^{-1}\big)\Gfive .
\end{align}
The commutator, with $D=D_+-D_-Y$ and $[D_+,D_-]=0$:
\begin{align}
[D,D_-]&=[D_+,D_-]-[D_-Y,D_-]\\
&=-\big(D_-YD_--D_-D_-Y\big)=-D_-[Y,D_-]\\
&=-D_-[Y,cW-1]=-c\,D_-[Y,W] .
\end{align}
The factor $D_-$ cancels against $D_-^{-1}$:
\begin{equation}
\boxed{\;D\dg\,\Gfive D^{-1}\Gfive=\Gfive\,(1-N)\,\Gfive,\qquad N=c\,[Y,W]\,D^{-1}\;}
\end{equation}
$[Y,W]=-\frac12[Y,K]$ is bounded (one hop in four dimensions, one in the fifth), so $N$ is large
only on the low modes of $D$. There the commutator flips chirality and moves the mode to the other
wall, out of the low-mode span, so $N^2$ is small and $(1-N)^{-1}\approx1+N$: GCR needs only a few
more steps than for the identity. Measured: 2--3 iterations more than the solve with $D$.

\subsection{On the even sites, step by step}
The solves of HMC are on the even sites. Since $\Gfive$ acts within a site it commutes with $\pi$
and $\iota$, and with Eq.~\eqref{eq:Sinv}
\begin{equation}
S^{-\dagger}=\pi D^{-\dagger}\iota\approx\pi\Gfive D^{-1}\Gfive\iota=\Gfive\,\pi D^{-1}\iota\,\Gfive
=\Gfive S^{-1}\Gfive .
\end{equation}
The solve of $D\dg w=r$:
\begin{align}
r'_e&=r_e-D_{oe}\dg(\Co\dg)^{-1}r_o &&\text{so that }S\dg w_e=r'_e,\\
w_e&=(\Ce\dg)^{-1}z &&\Longrightarrow\ S\dg(\Ce\dg)^{-1}z=r'_e\ \Longleftrightarrow\ K\dg z=r'_e,\\
z&=\text{GCR on }K\dg\text{ with preconditioner }\Ce\dg\,\Gfive S^{-1}\Gfive,\\
S^{-1}t&\approx\Ce^{-1}\,\mathcal M\,(\Spv\Cpve^{-1})\dg t &&\text{(one cycle of }\AR),\\
w_e&=(\Ce\dg)^{-1}z,\qquad w_o=(\Co\dg)^{-1}\big(r_o-D_{eo}\dg w_e\big).
\end{align}
The line $K\dg z=r'_e$ uses $K=\Ce^{-1}S$: $K\dg=S\dg(\Ce\dg)^{-1}$. The approximation for $S^{-1}$
comes from Eq.~\eqref{eq:ARinv}: $S^{-1}=\Ce^{-1}\AR^{-1}(\Spv\Cpve^{-1})\dg$. In the
preconditioner, $(\Spv\Cpve^{-1})\dg t$ is the right-hand side of $\AR$ for the source $(t,0)$ (no
odd part), so it is applied directly as one adjoint Schur factor of $\Dpv$. Each GCR step costs
one cycle, as for $D$.

% ======================================================================
\section{HMC}

\subsection{The even-odd Ratio action}
Two flavours over two Pauli--Villars flavours, as Bridge++ writes it (\code{Action\_F\_alt\_Ratio\_eo}):
with $K=\Ce^{-1}S$ (the left Schur factor) and $H=\Gfive K$,
\begin{equation}
S_F=\phi_e\dg H_{\rm PV}\,(H\dg H)^{-1}H_{\rm PV}\dg\phi_e .
\end{equation}
Since $\Gfive$ is Hermitian and unitary,
\begin{equation}
H\dg H=K\dg\Gfive\Gfive K=K\dg K ,
\end{equation}
and with $v_1\equiv H_{\rm PV}\dg\phi_e=K_{\rm PV}\dg\Gfive\phi_e$,
\begin{equation}
S_F=v_1\dg\,(K\dg K)^{-1}\,v_1 .
\end{equation}
\paragraph{Heat bath.} Draw $\xi_e$ Gaussian and set
$\phi_e=H_{\rm PV}(H_{\rm PV}\dg H_{\rm PV})^{-1}H\dg\xi_e$. Then, step by step,
\begin{align}
H_{\rm PV}\dg\phi_e&=H_{\rm PV}\dg H_{\rm PV}(H_{\rm PV}\dg H_{\rm PV})^{-1}H\dg\xi_e=H\dg\xi_e,\\
S_F&=(H\dg\xi_e)\dg(H\dg H)^{-1}(H\dg\xi_e)\\
&=\xi_e\dg H(H\dg H)^{-1}H\dg\xi_e\\
&=\xi_e\dg HH^{-1}H^{-\dagger}H\dg\xi_e=|\xi_e|^2 ,
\end{align}
as required. The heat bath costs one solve with the Pauli--Villars operator (CG; it is well
conditioned).

\paragraph{Force.} Vary the links. With $\chi=\Gfive\phi_e$, $v_1=K_{\rm PV}\dg\chi$ and
$v_2=(K\dg K)^{-1}v_1$:
\begin{align}
\delta S_F&=\delta v_1\dg\,v_2+v_2\dg\,\delta v_1+v_1\dg\,\delta\big[(K\dg K)^{-1}\big]v_1\\
&=2\,\Re\,(\delta v_1)\dg v_2-v_1\dg(K\dg K)^{-1}\,\delta(K\dg K)\,(K\dg K)^{-1}v_1\\
&=2\,\Re\,\chi\dg\,\delta K_{\rm PV}\,v_2-v_2\dg\big(\delta K\dg\,K+K\dg\,\delta K\big)v_2\\
&=2\,\Re\,\chi\dg\,\delta K_{\rm PV}\,v_2-2\,\Re\,(Kv_2)\dg\,\delta K\,v_2 .
\end{align}
So every force evaluation needs one solve $v_2=(K\dg K)^{-1}v_1$, and every energy evaluation (at
the start and the end of a trajectory) one more. The derivative of $K$ is that of the hops in $S$
($\Ce$ does not depend on the links), taken with respect to the smeared links and carried back
through the stout derivative.

\subsection{$(K\dg K)^{-1}$ with the multigrid: two solves}
Invert the product and use $K=\Ce^{-1}S$:
\begin{align}
(K\dg K)^{-1}&=K^{-1}K^{-\dagger}\\
&=(\Ce^{-1}S)^{-1}\big[(\Ce^{-1}S)\dg\big]^{-1}\\
&=S^{-1}\Ce\;\Ce\dg S^{-\dagger} .
\end{align}
So $z=(K\dg K)^{-1}r$ is computed as
\begin{align}
w&=\Ce\dg\,S^{-\dagger}r=\Ce\dg\,\pi D^{-\dagger}\iota\,r &&\text{(a solve with $D\dg$, Sec.~\ref{sec:dagger})},\\
z&=S^{-1}\Ce\,w=\pi D^{-1}\iota\,\Ce w &&\text{(a solve with $D$)}.
\end{align}
Each solve stops at a relative residual $10^{-10}$. The normal equation must be solved much more
accurately ($10^{-16}$ for the force, $10^{-24}$ for the energy), so the two solves are repeated on
the residual (iterative refinement):
\begin{align}
r^{(k)}&=b-K\dg Kz^{(k)}\\
z^{(k+1)}&=z^{(k)}+K^{-1}K^{-\dagger}r^{(k)}\qquad(\text{two solves to }10^{-10}),
\end{align}
each sweep gaining about ten digits, at most eleven sweeps.

\paragraph{When the sweeps stall.} If the sweeps stop gaining (the coarse space no longer fits the
links, as after large link changes), the solve is repeated from the current $z$ with $\Pro$ and the
coarse operators rebuilt for the current links, then also after a bootstrap; if it still stalls,
CG on $K\dg K$ finishes it from the current $z$ (with at least 50000 iterations), with a warning.
A trajectory is never lost to the multigrid.

\subsection{Neff's $\alpha$ in HMC}
\paragraph{CG.} From Eq.~\eqref{eq:Ksim}, $K=Q_eK^{(\alpha)}Q_e^{-1}$. Then
\begin{align}
K\dg K&=Q_e^{-1}K^{(\alpha)\dagger}Q_eQ_eK^{(\alpha)}Q_e^{-1}\\
&=Q_e^{-1}\big(Q_eK^{(\alpha)}\big)\dg\big(Q_eK^{(\alpha)}\big)Q_e^{-1}\\
&=Q_e^{-1}\,(KQ_e)\dg(KQ_e)\,Q_e^{-1}\qquad(KQ_e=Q_eK^{(\alpha)}),\\
(K\dg K)^{-1}&=Q_e\big[(KQ_e)\dg(KQ_e)\big]^{-1}Q_e .
\end{align}
So CG runs on the normal equation of $KQ_e$, whose singular values are close to those of
$\alpha K^{(\alpha)}$, and the action and the force stay those of $K$. On $24^4$ at $m=0.001$ this
cut the CG iterations by 22\%.

\paragraph{Multigrid.} The solve with $D$ uses $D^{-1}=QD^{(\alpha)\,-1}$: the left form
$\AL^{(\alpha)}$ with the factors of the table above and $\psi=Q\psi^{(\alpha)}$. The solve with $D\dg$ is
unchanged: $D^{(\alpha)\dagger}=QD\dg$, so $\alpha$ would only rescale its right-hand side. The setup
is that of $D$.

\paragraph{Why $\alpha$ helps CG more than the multigrid.} $K^{(\alpha)}=Q_e^{-1}KQ_e$ has the
eigenvalues of $K$ but other singular values ($Q_e$ is not unitary). CG on the normal equation sees
the singular values, so it gains. The multigrid's iterations do not follow the whole condition
number: the coarse grid takes the low modes and the smoother the rest, and $\alpha$ helps only the
smoother's part. In the propagator mass scan this saves 2--3 iterations at every mass with
$M_5=1.8$ (smoother interval $[\lambda_{\rm hi}/30,\lambda_{\rm hi}]$, wide) and only at heavy
masses with stout links ($[\lambda_{\rm hi}/10,\lambda_{\rm hi}]$).

\subsection{What is kept along a trajectory}
\begin{itemize}
\item The test vectors $y$: kept from solve to solve (on the host between solves), improved by a
bootstrap (from the old $y$, not a new filter) when the pseudofermions are drawn.
\item $\Pro$, the coarse operators, the top of the smoother's interval: kept from solve to solve
while the links move, and rebuilt for the current links when the GCR iterations per solve have
grown by half (\code{rebuild\_iteration\_growth}~$=1.5$) or after a bootstrap.
\item The cycle itself is made for each solve from the kept parts and dropped after it.
\end{itemize}

\subsection{Memory (24$^4$, $m=0.001$, 20 vectors, levels below, A30)}
\begin{center}
\begin{tabular}{@{}lr@{}}
\toprule
part & GB\\
\midrule
HMC itself (links, smeared levels, momenta, pseudofermion, operators) & 3.3\\
multigrid kept between solves ($\Pro$ 4.75) & 6.6\\
fields of the sweeps & 2.6\\
the solve with $D\dg$ on top (GCR, Schur factors, cycle) & 6.4\\
\midrule
peak (live) & 19.1\\
\bottomrule
\end{tabular}
\end{center}

\subsection{What the measurements say so far}
\begin{itemize}
\item On a thermalized dynamical $16^4$ configuration the multigrid reproduces the trajectories
of CG (plaquette to $10^{-11}$).
\item There CG needs about 300 iterations per solve at $m=0.01$ and at $m=0.001$ alike: the low end
of the spectrum is held up by the finite volume, not by the mass. The multigrid needs about 55
iterations, but one of them (a whole cycle) costs about 24 CG iterations, so it is 3--4 times
slower per trajectory. Tuning gains at most 6\%.
\item On the quenched $24^4$ configuration CG needs about 6400 iterations and the multigrid wins,
but a trajectory of length 1 there blows up. So there are no data yet that show the multigrid
faster in HMC; it needs a volume where CG takes more than about 1300 iterations per solve.
\end{itemize}

% ======================================================================
\section{Map to the code}
\begin{center}
\small
\begin{tabular}{@{}p{0.36\linewidth}p{0.6\linewidth}@{}}
\toprule
step & where\\
\midrule
$C$, $C^{-1}$, hops, $\Gfive$, $Q$ on the device & \code{src/backend/cuda\_bridge/mdwf.rs};\newline \code{cuda\_bridge/include/fermion/}\\
even-odd forms, $\AR$, $\AL$ (\code{PvRatio}, \code{SchurFactor}, \code{AfterQ}) & \code{src/fermion/even\_odd.rs}\\
GCR, FGMRES, BiCGStab, MINRES, CG & \code{src/solver/}\\
Chebyshev smoother, two-grid cycle & \code{src/solver/preconditioner/multigrid/cycle.rs};\newline \code{src/backend/cuda\_bridge/multigrid.rs}\\
factorized correction & \code{src/solver/preconditioner/multigrid/correction.rs};\newline \code{src/backend/cuda\_bridge/correction.rs}\\
$\Pro$: fold, restrict, prolong, CholeskyQR2 & \code{src/backend/cuda\_bridge/coarse.rs};\newline \code{cuda\_bridge/include/multigrid/fold.cuh}\\
coarse operators by probing & \code{CudaMdwf::coarse\_blocks} in\newline \code{src/backend/cuda\_bridge/mdwf.rs}\\
levels below & \code{src/backend/cuda\_bridge/levels.rs}\\
filter, bootstrap & \code{low\_modes} in\newline \code{src/solver/preconditioner/multigrid/mdwf\_even\_odd.rs};\newline \code{bootstrap} in \code{src/backend/cuda\_bridge/multigrid.rs}\\
solve with $D\dg$ & \code{adjoint\_by\_gamma5} in \code{src/fermion/even\_odd.rs}\\
HMC: two solves, sweeps, retries & \code{src/hmc/dwf\_multigrid.rs};\newline \code{src/backend/cuda\_bridge/hmc\_multigrid.rs}\\
HMC: action, force, CG with $\alpha$, CG fallback & \code{src/hmc/even\_odd.rs}; \code{src/hmc/solve.rs}\\
\bottomrule
\end{tabular}
\end{center}

\end{document}
